%%%%%
%
\section{Indice de Fogging des images d'une sonde PMS 2D-C et nettoyage des donn\'ees}

\itemstyle{Op\'erateur~:} {\sf Fogging\_2D}
\index{Fogging\_2D}

%%%%%OBJET%%%%
\itemstyle{Objet~:}
Cet op�rateur calcule un indice de Fogging \`a partir des donn\'ees images 
d'une sonde PMS 2D-C. Un nettoyage des donn\'ees est effectu\'e
optionnellement.

%%%%%ARGUMENTS%%%%
\bigskip
%%%%%ENTREES%%%%
\begin{tabarg}
\multicolumn{3}{l}{\titretab {Entr\'ees}}\\
\hline
$Time_{reg}$ & Vecteur, {\tt Float\_64} & Temps avec fr\'equence de 1 Hertz ($s$)\\
$Image_{2D}$ & Matrice X1024, {\tt Int\_32} & Image\\
$Time_{asyn}$ & Vecteur, {\tt Float\_64} & Temps de la variable d'entr\'ee Image ($s$)\\
$TAS_{2D}$ & Vecteur, {\tt Int\_16} & True Air Speed ($m/s$)\\
$Code\_Hexa\_sonde$ & Coefficient, {\tt String} & Code hexadecimal de la sonde\\
$Masque\_sonde$ & Coefficient, {\tt String} & Masque hexadecimal des donn\'ees images\\
$Bits\_temps$ & Coefficient, {\tt Int\_16} & Nombre de bits de codage du temps dans les s\'eparateurs\\
$Nbval_{as}$ & Coefficient, {\tt Int\_32} & Nombre d'\'el\'ements dans le vecteur de temps ascynchrone\\
$Temps_{debut}$ & Coefficient, {\tt Int\_32} & Temps du d\'ebut de l'analyse en secondes depuis minuit\\
$Temps_{fin}$ & Coefficient, {\tt Int\_32} & Temps de la fin de l'analyse en secondes depuis minuit\\
$Tronq_{debut}$ & Coefficient, {\tt Int\_16} & Num\'ero de la diode \`a consid\'erer (d\'ebut)\\
$Tronq_{fin}$ & Coefficient, {\tt Int\_16} & Num\'ero de la diode \`a consid\'erer (fin)\\
$Freq$ & Coefficient, {\tt Int\_16} & Nombre d'enregistrements \`a traiter ensemble pour la d\'etection de Fogging\\
$Nettoyer$ & Coefficient, {\tt Int\_16} & Si on doit nettoyer les donn\'ees\\
\end{tabarg}

\medskip
%%%%%SORTIES%%%%
\begin{tabarg}
\multicolumn{3}{l}{\titretab {Sorties}}\\
\hline
$Fogging_{2D}$ & Vecteur, {\tt Int\_32} & Indice de Fogging\\
$Images\_nettoyees_{2D}$ & Matrice X1024, {\tt Int\_32} & Donn\'ees Image de la sonde 2D nettoy\'ees si demand\'e\\
\end{tabarg}


%%%%%FORMULE%%%%
\itemstyle{Formule}

%%% Fogging_2D


%%%%%COMMENTAIRE%%%%
\itemstyle{Commentaire~:}
Les donn�es images de la sonde PMS-2D-C sont analys�es par imagette.
La m�thode de d�tection des imagettes a �t� modifi�e afin d'utiliser la m�me
que celle de l'op�rateur Compte\_img\_2D.
Les s�parateurs, lorsqu'il y a fogging des donn�es, peuvent �tre
et sont souvent endommag�s. Ceci fausse beaucoup les calculs de spectre, de
compte de particules et de d�tection de types de particules. La seule m�thode
de calcul qui est la moins touch�e est le calcul des spectres par la m�thode
� Bill Hall (NCAR) car celle-ci fait non pas de l'analyse d'imagettes, mais
plut�t de l'analyse de groupes de pixels. Le seul impact dans ce cas-ci est
la pr�sence de bruit dans les donn�es, m�me apr�s nettoyage, ou la
sous-estimation de la taille caus�e par le nettoyage et la suppression de
pixels et la moins bonne pr�cision dans le positionnement des particules dans
le temps.
La m�thode utilis�e pour calculer l'indice de Fogging est de compter le nombre
d'occurences des patron 010 et 0110 dans les colonnes de donn�es images.
L'option de nettoyage des donn�es remplace les patrons 010 et 0110 par 000 et
0000 respectivement.
Pour des donn�es bruit�es par le fogging, l'ordre de grandeur de l'indice se
situe aux environs de 10\^4. Une plage d'enregistrements ayant un indice de
fogging moins de 800 est consid�r� comme n'ayant pas de fogging.

%%%%%AUTEUR%%%%
\itemstyle{Auteur~:}
CNRM/GMEI/MNP

%%%%%REFERENCE%%%%
\itemstyle{R\'ef\'erences~:}
